\documentclass[11pt,a4paper]{article}

%========================
% MOTEUR & FRANÇAIS
%========================
\usepackage{fontspec}         % LuaLaTeX/XeLaTeX
\usepackage[french]{babel}
\usepackage[final]{microtype} % microtypographie (jolie justification)
\usepackage{siunitx}          % nombres FR (espaces fines, virgule)
\sisetup{locale=FR, group-minimum-digits=4}

%========================
% MATHS & MISE EN PAGE
%========================
\usepackage{amsmath,amssymb}
\usepackage{geometry}
\geometry{top=1.5cm,bottom=2cm,left=1.8cm,right=1.8cm}
\usepackage{fancyhdr}
\setlength{\headheight}{15pt}
\usepackage{enumitem}
\usepackage{multicol}
\usepackage{tikz}
\usepackage{pgfplots}
\pgfplotsset{compat=1.18}
\usepackage{xcolor}
\usepackage{graphicx}
\usepackage{array}
\usepackage{fontawesome5}
\usepackage{tabularx}
\usepackage[most]{tcolorbox}
\tcbuselibrary{skins,breakable}

%========================
% TOGGLE CORRIGÉ
%========================
\newif\ifcorrige
\corrigefalse          % <-- Version Élève (par défaut)
%\corrigetrue         % <-- Décommenter pour la Version Corrigé

% Réponse (s'affiche seulement en mode corrigé)
\definecolor{warning}{RGB}{231,76,60}
\DeclareRobustCommand{\rep}[1]{%
	\ifcorrige
	\begingroup\color{warning}\ifmmode #1\else \textbf{#1}\fi\endgroup
	\fi
}
% Afficher des blocs entiers seulement en mode corrigé
\newcommand{\onlycorrige}[1]{\ifcorrige #1\fi}

%========================
% PALETTE COULEURS
%========================
\definecolor{primary}{RGB}{41,128,185}      % Bleu
\definecolor{secondary}{RGB}{52,73,94}      % Gris foncé
\definecolor{accent}{RGB}{46,204,113}       % Vert
\definecolor{light}{RGB}{236,240,241}       % Gris clair
\definecolor{bglight}{RGB}{247,249,250}     % Fond clair

%========================
% EN-TÊTE & PIED DE PAGE
%========================
\pagestyle{fancy}
\fancyhf{}
\fancyhead[L]{\color{secondary}\textbf{6\textsuperscript{e}} \,|\, Mathématiques}
\fancyhead[R]{\color{secondary}\textbf{\ifcorrige Corrigé — Nombres Entiers\else Nombres Entiers\fi}}
\fancyfoot[C]{\color{secondary}\thepage}
\renewcommand{\headrulewidth}{0pt}
\renewcommand{\footrulewidth}{0pt}

%========================
% BOÎTES
%========================
\newtcolorbox{exercisebox}[2][]{
	colback=white,colframe=primary,
	fonttitle=\bfseries\large,
	title={#2},
	arc=2mm,boxrule=1.5pt,
	left=15pt,right=15pt,top=12pt,bottom=12pt,
	boxsep=0pt,breakable,enhanced,
	width=\linewidth,
	grow to left by=3mm, % 3mm = plus sûr contre overfull
	grow to right by=3mm,
	#1
}
\newtcolorbox{consignebox}{
	colback=light!30,colframe=secondary,
	arc=3mm,boxrule=0pt,leftrule=4pt,
	left=15pt,right=15pt,top=12pt,bottom=12pt,
	enhanced,width=\linewidth,
	grow to left by=3mm,
	grow to right by=3mm
}

%========================
% MACROS UTILES
%========================
\newcommand{\points}[1]{\hfill \colorbox{primary}{\color{white}\textbf{\,#1 pts\,}}}
\newcommand{\badge}[2]{%
	\tikz[baseline=(char.base)]{%
		\node[shape=rectangle,draw,thick,rounded corners=2mm,inner sep=2pt,fill=#1,text=white,font=\footnotesize\bfseries] (char) {#2};}%
}
\newcommand{\dotline}[1]{\makebox[#1][s]{\dotfill}}
\newcolumntype{L}[1]{>{\hsize=#1\hsize\arraybackslash}X}
\newcommand{\zonereponse}[1]{%
	\par\noindent\colorbox{light}{\parbox{\dimexpr\linewidth-2\fboxsep\relax}{\vspace{#1}}}\par}
\newcommand{\zonereponseLignes}[2][0.9cm]{%
	\par\noindent\colorbox{light}{\parbox{\dimexpr\linewidth-2\fboxsep\relax}{%
			\setlength{\parskip}{#1}\foreach \i in {1,...,#2}{\noindent\dotfill\par}}}\par}

%========================
\begin{document}
	%========================
	
	% Filigrane CORRIGÉ (optionnel)
	\onlycorrige{
		\begin{tikzpicture}[remember picture,overlay]
			\node[rotate=30,scale=8,text=secondary!12] at (current page.center) {CORRIGÉ};
		\end{tikzpicture}
	}
	
	% Bandeau supérieur
	\begin{tikzpicture}[remember picture,overlay]
		\fill[primary] (current page.north west) rectangle ([yshift=-3.5cm]current page.north east);
		\node[anchor=north west,text=white,font=\Huge\bfseries]
		at ([xshift=2cm,yshift=-1cm]current page.north west)
		{\ifcorrige Corrigé — Contrôle n° 1 \else Contrôle n° 1\fi};
		\node[anchor=north west,text=white,font=\Large]
		at ([xshift=2cm,yshift=-2cm]current page.north west) {Les Nombres Entiers};
		\node[anchor=north east,text=white,font=\large]
		at ([xshift=-2cm,yshift=-1.5cm]current page.north east) {\faIcon{clock} 50 minutes};
		\node[anchor=north east,text=white,font=\large]
		at ([xshift=-2cm,yshift=-2.3cm]current page.north east) {\faIcon{calculator} Calculatrice interdite};
	\end{tikzpicture}
	
	\vspace*{1.5cm}
	
%========================
% CARTOUCHE ÉLÈVE (NOM large)
%========================
\begin{tcolorbox}[
	colback=bglight,colframe=secondary,arc=3mm,boxrule=1pt,
	left=15pt,right=15pt,top=10pt,bottom=10pt,
	enhanced,width=\linewidth,
	grow to left by=3mm, grow to right by=3mm
	]
	\noindent
	\begin{minipage}{0.78\linewidth}
		\setlength{\tabcolsep}{4pt}
		
		% --- NOM (1 ligne entière) ---
		\begin{tabularx}{\linewidth}{@{} l L{2.8} @{}}
			\faIcon{user}\ \textbf{NOM :} & \dotfill \\
		\end{tabularx}
		
		\vspace{0.35cm}
		
		% --- Prénom + Classe ---
		\begin{tabularx}{\linewidth}{@{} l L{0.8} l L{0.5} @{}}
			\textbf{Prénom :} & \dotfill &
			\textbf{Classe :} & \dotfill \\
		\end{tabularx}
		
		\vspace{0.35cm}
		
		% --- Date (pleine largeur) ---
		\begin{tabularx}{\linewidth}{@{} l L{1.2} @{}}
			\faIcon{calendar}\ \textbf{Date :} & \dotfill \\
		\end{tabularx}
	\end{minipage}\hfill
	\begin{minipage}{0.20\linewidth}
		\raggedleft
		\fcolorbox{secondary}{white}{\textbf{NOTE : }\makebox[1cm]{\rule{0pt}{1.1em}}\textbf{ / 20}}
	\end{minipage}
\end{tcolorbox}

	
	%========================
	% BARÈME
	%========================
	\vspace{0.2cm}
	\begin{center}
		\badge{primary}{Ex.1: 4pts}\quad
		\badge{primary}{Ex.2: 5pts}\quad
		\badge{primary}{Ex.3: 4pts}\quad %\\[-2pt]
		\badge{primary}{Ex.4: 4pts}\quad
		\badge{primary}{Ex.5: 3pts}
		%\badge{accent}{Bonus: +2pts}
	\end{center}
	
	%========================
	% CONSIGNES
	%========================
	\vspace{0.2cm}
	\begin{consignebox}
		\faIcon{info-circle} \textbf{CONSIGNES IMPORTANTES}
		\begin{itemize}[leftmargin=*,noitemsep,topsep=5pt]
			\item[\faIcon{check}] Écrivez lisiblement et justifiez vos réponses
			\item[\faIcon{check}] Respectez l'orthographe des nombres en lettres
			\item[\faIcon{check}] Les exercices peuvent être traités dans n'importe quel ordre
			\item[\faIcon{times}] Calculatrice non autorisée
		\end{itemize}
	\end{consignebox}
	
	%========================
	% EXERCICE 1
	%========================
	\vspace{0.5cm}
	\begin{exercisebox}{\faIcon{sort-numeric-down} Exercice 1 : Rang des chiffres \points{4}}
		Considérons le nombre : \textbf{\num{7256394}}
		
		\begin{enumerate}[label=\textcolor{primary}{\textbf{\alph*)}},leftmargin=*]
			\item Quel est le chiffre des centaines ? \dotfill \points{1}
			\item Quel est le chiffre des dizaines de milliers ? \dotfill \points{1}
			\item Quel est le \emph{nombre de milliers} dans ce nombre ? \dotfill \points{1}
			\item Donne la décomposition décimale de ce nombre : \points{1}
			\zonereponseLignes[0.5cm]{2}
		\end{enumerate}
	\end{exercisebox}
	
	%========================
	% EXERCICE 2
	%========================
	\begin{exercisebox}{\faIcon{spell-check} Exercice 2 : Écriture en lettres \points{5}}
		\textbf{Écris les nombres suivants en toutes lettres :}
		
		\begin{enumerate}[label=\textcolor{primary}{\textbf{\alph*)}},leftmargin=*]
			\item \num{384} : \dotfill \points{1}
			\item \num{1200} : \dotfill \points{1}
			\item \num{5087} : \dotfill \par \dotfill \points{1}
			\item \num{14680} : \dotfill \par \dotfill \points{1}
			\item \num{300092} : \dotfill \par \dotfill \points{1}
		\end{enumerate}
	\end{exercisebox}
	
	%========================
	% EXERCICE 3
	%========================
	\begin{exercisebox}{\faIcon{ruler-horizontal} Exercice 3 : Demi-droite graduée \points{4}}
		\vspace{0.3cm}
		\textbf{Observe la demi-droite graduée ci-dessous :}
		
		\begin{center}
			\begin{tikzpicture}[scale=1]
				% Axe
				\draw[thick,->] (0,0) -- (12,0);
				% Graduations principales (tous les 10)
				\foreach \x in {0,2,4,6,8,10} {
					\draw[thick] (\x,0.1) -- (\x,-0.1);
					\node[below] at (\x,-0.3) {\ifnum\x=0 0\else \x0\fi};
				}
				% Graduations secondaires (tous les 5)
				\foreach \x in {1,3,5,7,9,11} {\draw (\x,0.05) -- (\x,-0.05);}
				% Points
				\node[circle,fill=primary,inner sep=2pt,label=above:$A$] at (3,0) {};
				\node[circle,fill=primary,inner sep=2pt,label=above:$B$] at (7,0) {};
				\node[circle,fill=primary,inner sep=2pt,label=above:$C$] at (9.5,0) {};
			\end{tikzpicture}
		\end{center}
		
		\begin{enumerate}[label=\textcolor{primary}{\textbf{\alph*)}},leftmargin=*]
			\item Quelle est l'abscisse du point $A$ ? \dotfill \points{1}
			\item Quelle est l'abscisse du point $B$ ? \dotfill \points{1}
			\item Quelle est l'abscisse du point $C$ ? \dotfill \points{1}
			\item Place le point $D$ d'abscisse \num{45} sur la demi-droite ci-dessus. \points{1}
		\end{enumerate}
	\end{exercisebox}
	
	%========================
	% EXERCICE 4
	%========================
	\begin{exercisebox}{\faIcon{greater-than-equal} Exercice 4 : Comparaison et rangement \points{4}}
		\begin{enumerate}[label=\textcolor{primary}{\textbf{\alph*)}},leftmargin=*]
			\item Compare avec $<$, $>$ ou $=$ : \points{2}
			\begin{multicols}{4}
				\begin{itemize}[label=$\bullet$]
					\item \num{3456} \dotfill \num{3546}
					\item \num{12089} \dotfill \num{12098}
					\item \num{7000} \dotfill \num{6999}
					\item \num{45230} \dotfill \num{45203}
				\end{itemize}
			\end{multicols}
			\item Range dans l'ordre croissant : \points{2}
			\begin{center}
				\num{8743}\,;\quad \num{8347}\,;\quad \num{8734}\,;\quad \num{8437}\,;\quad \num{8374}
			\end{center}
			\zonereponseLignes[0.4cm]{2}
		\end{enumerate}
	\end{exercisebox}
	
	%========================
	% EXERCICE 5
	%========================
	\begin{exercisebox}{\faIcon{lightbulb} Exercice 5 : Problème \points{3}}
		\begin{tcolorbox}[colback=blue!5,colframe=primary!30,arc=2mm,boxrule=0.5pt]
			\textbf{Contexte :} La bibliothèque du collège compte ses livres.
			\begin{itemize}[noitemsep,topsep=5pt]
				\item Il y a \textbf{3 milliers} de romans
				\item Il y a \textbf{15 centaines} de livres documentaires
				\item Il y a \textbf{85 dizaines} de bandes dessinées
			\end{itemize}
		\end{tcolorbox}
		
		\begin{enumerate}[leftmargin=*]
			\item \textbf{Calcule} le nombre total de livres : \points{2}
			\zonereponseLignes[0.5cm]{3}
			\item \textbf{Écris} ce nombre total en toutes lettres : \points{1}
			\zonereponseLignes[0.5cm]{2}
		\end{enumerate}
	\end{exercisebox}
	
	%========================
	% AUTO-ÉVAL
	%========================
	\vspace{0.3cm}
	\begin{center}
		\begin{tikzpicture}
			\node[draw=secondary,rounded corners=3mm,inner sep=8pt,fill=light!30] {
				\begin{tabular}{lccccc}
					\textbf{Auto-évaluation :} &
					\faIcon{frown} Difficile &
					\faIcon{meh} Moyen &
					\faIcon{smile} Facile &
					\phantom{xx} &
					\textbf{Temps utilisé :} \underline{\phantom{xxxxx}} min
				\end{tabular}
			};
		\end{tikzpicture}
	\end{center}
	
	%========================
	% CORRIGÉ (visible seulement si \corrigetrue)
	%========================
	\onlycorrige{
		\newpage
		\begin{center}\Large\textbf{CORRIGÉ}\end{center}
		
		\begin{exercisebox}{\faIcon{sort-numeric-down} Exercice 1 – Corrigé}
			Pour le nombre \num{7256394} :
			\begin{enumerate}[label=\textcolor{primary}{\textbf{\alph*)}},leftmargin=*]
				\item Chiffre des centaines : \rep{3}
				\item Chiffre des dizaines de milliers : \rep{5}
				\item Nombre de milliers : \rep{\num{7256}} \; (car \num{7256394} = \num{7256}\,$\times$\,\num{1000} + \num{394})
				\item Décomposition : \rep{$7\times\num{1000000} + 2\times\num{100000} + 5\times\num{10000} + 6\times\num{1000} + 3\times\num{100} + 9\times\num{10} + 4$}
			\end{enumerate}
		\end{exercisebox}
		
		\begin{exercisebox}{\faIcon{spell-check} Exercice 2 – Corrigé}
			\begin{enumerate}[label=\textcolor{primary}{\textbf{\alph*)}},leftmargin=*]
				\item \num{384} : \rep{trois-cent-quatre-vingt-quatre}
				\item \num{1200} : \rep{mille-deux-cents}
				\item \num{5087} : \rep{cinq-mille-quatre-vingt-sept}
				\item \num{14680} : \rep{quatorze-mille-six-cent-quatre-vingts}
				\item \num{300092} : \rep{trois-cent-mille-quatre-vingt-douze}
			\end{enumerate}
		\end{exercisebox}
		
		\begin{exercisebox}{\faIcon{ruler-horizontal} Exercice 3 – Corrigé}
			\begin{enumerate}[label=\textcolor{primary}{\textbf{\alph*)}},leftmargin=*]
				\item Abscisse de $A$ : \rep{\num{30}}
				\item Abscisse de $B$ : \rep{\num{70}}
				\item Abscisse de $C$ : \rep{\num{95}}
				\item $D$ à l’abscisse \num{45} (au milieu entre \num{40} et \num{50} sur la figure). \rep{Placer $D$ à $x=4.5$ sur l’axe}
			\end{enumerate}
		\end{exercisebox}
		
		\begin{exercisebox}{\faIcon{greater-than-equal} Exercice 4 – Corrigé}
			\textbf{a)} \quad
			\rep{\num{3456} < \num{3546}},\quad
			\rep{\num{12089} < \num{12098}},\quad
			\rep{\num{7000} > \num{6999}},\quad
			\rep{\num{45230} > \num{45203}}.
			
			\medskip
			\textbf{b)} Ordre croissant : \rep{\num{8347} < \num{8374} < \num{8437} < \num{8734} < \num{8743}}.
		\end{exercisebox}
		
		\begin{exercisebox}{\faIcon{lightbulb} Exercice 5 – Corrigé}
			\textbf{Calculs} : \quad
			Romans : \rep{\num{3000}} ;\;
			Documentaires : \rep{\num{1500}} ;\;
			BD : \rep{\num{850}}. \\
			Total : \rep{\num{5350}} \, livres. \quad
			En lettres : \rep{cinq-mille-trois-cent-cinquante}.
		\end{exercisebox}
		
		\begin{tcolorbox}[colback=accent!10,colframe=accent]
			\textbf{BONUS – Corrigé}
			
			Le nombre a la forme \(\overline{5c3u}\) avec \(c=2u\) et \(5+c+3+u=16\).
			Donc \(8+3u=16 \Rightarrow u=\rep{2}\) puis \(c=\rep{4}\).  
			Le nombre est \rep{\num{5232}} (ou \rep{\num{5632}} selon l’interprétation « centaines = double des unités »).  
			\textit{NB : Dans l’énoncé donné, pour obtenir exactement 16 avec « centaines = double des unités » et « dizaines = 3 », la solution cohérente est \(\overline{5\,6\,3\,2}\) : \(5+6+3+2=16\).}
		\end{tcolorbox}
	}
	%========================
\end{document}
